
A central assumption underlying the use of adversarial NLP benchmarks is that harder, adversarially constructed inputs expose model failures in a way that informs deployment risk. Much of the resulting literature evaluates accuracy degradation: a model's classification accuracy drops from a clean-data baseline to an adversarial one. This work asks a different question: under adversarial conditions, does model confidence remain informative about correctness?

Calibration---the alignment between predicted confidence and observed accuracy---is a critical reliability property. A model that assigns 90\% confidence to its predictions should be correct approximately 90\% of the time; deviations from this alignment render confidence signals unreliable as a safety or quality indicator. Expected Calibration Error (ECE) \citep{Guo2017Calibration} is the standard metric for measuring this alignment. Minderer et al. [2021] demonstrated that distribution shift in vision models consistently degrades ECE; the analogous question for adversarial distribution shift in NLP has not been directly measured with logit-based ECE for open-weight LLMs.

This paper fills this gap. The central finding is counterintuitive: evaluating Llama-2-7b-hf on ANLI R1---a benchmark specifically constructed so that human annotators could fool state-of-the-art NLP models---produces a calibration \textit{improvement} (ECE = 0.239 versus the clean-data baseline ECE = 0.279, $\Delta ECE = -0.041)$. The benchmark designed to expose model failures made the model appear better calibrated. This outcome is not an error; it reveals a structural property of model-in-the-loop adversarial benchmark construction.

The key insight is that model-in-the-loop adversarial examples (ANLI) are selected precisely because the model fails on them. If the model fails with appropriately low confidence, it is, by definition, well-calibrated on those failures. Human-adversarial examples (AdvGLUE MNLI), by contrast, are designed to exploit surface features to trigger high model confidence on incorrect answers---the canonical miscalibration mechanism. The distinction between these construction methods determines whether calibration degrades or improves.

\subsubsection{1.1 Problem Statement}

No published work measures logit-based ECE on adversarial NLP benchmark splits for open-weight LLMs, or examines how adversarial construction method and RLHF alignment jointly determine calibration outcomes. Adversarial robustness evaluation (accuracy-only) and calibration measurement (clean-data-only) have developed as separate research communities. This work occupies the intersection.

\subsubsection{1.2 Research Questions}

This work addresses four research questions:

\begin{itemize}
\item \textbf{RQ1 (Existence):} Does adversarial perturbation measurably increase ECE for open-weight LLMs on NLP benchmarks?
\item \textbf{RQ2 (Validity):} Is $\Delta ECE$ a valid calibration signal, or does it reflect label noise introduced by adversarial perturbation?
\item \textbf{RQ3 (Conditionality):} Is calibration degradation universal across task types and adversarial construction methods?
\item \textbf{RQ4 (RLHF Moderation):} Does RLHF alignment moderate adversarial calibration degradation, and does the effect depend on benchmark construction method?
\end{itemize}

\subsubsection{1.3 Contributions}

Three findings are reported:

\begin{enumerate}
\item \textbf{First measurement of logit-based ECE on adversarial NLP benchmarks.} $\Delta ECE$ = +0.071 for AdvGLUE MNLI confirms that adversarial calibration degradation is real and measurable for an open-weight LLM. Label preservation rate = 1.000 for all adversarial splits validates $\Delta ECE$ as a genuine calibration signal. The JSONL logit caching strategy employed here enables multiple downstream calibration analyses from a single inference pass.
\end{enumerate}

\begin{enumerate}
\item \textbf{Evidence that calibration degradation is task-type and construction-method conditional.} Human-adversarial NLI examples (AdvGLUE MNLI: $\Delta ECE$ = +0.071) degrade calibration; model-in-the-loop adversarial NLI examples (ANLI R1/R2: $\Delta ECE$ < 0) improve it. The ANLI difficulty gradient (R1 $\rightarrow$ R2 $\rightarrow$ R3) shows a monotonic transition from improvement to degradation. Binary classification consistently shows reversed $\Delta ECE$.
\end{enumerate}

\begin{enumerate}
\item \textbf{Evidence that RLHF alignment conditionally moderates adversarial calibration degradation.} Confirmed for model-in-the-loop adversarial (ANLI: 100\% moderation rate, 3/3 rounds); reversed for static human-adversarial benchmarks (AdvGLUE: alignment tax, $\Delta \Delta ECE = -0.026$, chat shows higher $\Delta ECE$ than base). This is the first documented benchmark-type $\times$ RLHF calibration interaction.
\end{enumerate}

All results are from a single-model pilot study. Generalization across model families remains an open question.

